\section{Introduction}

The rotational dynamics of thin rigid disks (popularly know as Euler's disks) has been a subject of study for many centuries (Euler, 18th Century). It is a phenomena of common experience that if a circular disk--such a coin--is spun upon a table, it then comes to rest abruptly with a sound of increasing frequency. As the disk rolls on its edge, the point of contact $P$ describes a circle with angular velocity $\Omega$. In an ideal, non-dissipative system, $\Omega$ is constant and the motion of the disk never stops. In real systems, the energy of a system is gradually dissipated until the object in motion comes to rest; in the case of an Euler disk under rotation in a plane, the precession frequency of the object gradually approaches infinity as the energy is dissipated. Although the analytical equations of motion of this system are well known, the determination of the dominant dissipative mechanism has been in debate in recent times: Several articles were published in the early 2000s \cite{Moffatt2000} debating the nature of the main dissipative component of the system \cite{PhysRevE2002, vandenEngh2000, Petrie2002}, and from such debates was established the consensus that the main component of the dissipation forces is the friction with the plane. An ArXiV preprint \cite{Thorne2026} from March 2026 challenges this view for the very high precession frequency regime and bringing this exciting phenomena once more to the spotlight. Most studies on the subject rely on computational methods or complex video camera setups in vacuum for evaluation, but spectral analysis techniques have been proven sufficiently efficient at accurately measuring its oscillation dynamics \cite{Darmendrail2021}. Therefore, we would like to propose an experimental spectral analysis of the rotational dynamics of an Euler disk using inexpensive equipment. We anticipate that such an experiment could be proven useful for analysing the dominant mechanism for the energy dissipation of the system at the finite-time singularity region.